\section{Continuum and relational geometry}

\begin{frame}{From blocks to kernels}
  \bookref{Chapter 3, Section 2}
  A $k$-block density matrix determines a step function
  \[
    \rho_k:[0,1]^2\to[0,1].
  \]
  As the partition is refined, this motivates a measurable continuum profile
  and a functional
  \[
    \mathcal F_\beta[\rho]
      =\beta\Phi[\rho]
       -\int_0^1\!\int_0^1 h(\rho(x,y))\,dx\,dy.
  \]

  \begin{definition}
    An undirected graphon is a symmetric measurable
    $W:[0,1]^2\to[0,1]$ \cite{lovaszszegedy2006}.
  \end{definition}

  Directed relational kernels need not satisfy $W(x,y)=W(y,x)$. Dense graph
  limits therefore supply a rigorous neighboring theory, but not every kernel
  model is an undirected graphon model.
\end{frame}

\begin{frame}{Continuous symmetry breaking}
  \bookref{Chapter 3, Sections 6--7}
  \small
  Consider a differentiable $O(m)$-invariant order-parameter field, $m\geq1$,
  \[
    \mathcal F[\phi]
      =\int_\Omega
      \left(\frac\kappa2|\nabla\phi|^2+a|\phi|^2+b|\phi|^4\right)dx,
      \qquad b,\kappa>0.
  \]
  For $a<0$, constant minima satisfy
  \[
    |\phi|=v=\sqrt{\frac{-a}{2b}},
    \qquad
    O(m)\longrightarrow O(m-1),
  \]
  and the vacuum manifold is
  \[
    O(m)/O(m-1)\cong S^{m-1}.
  \]

  In the constrained vacuum model with $m=2$, writing
  $\phi=v(\cos\theta,\sin\theta)$ gives
  \[
    \Delta\mathcal F=\frac{\kappa v^2}{2}\int_\Omega|\nabla\theta|^2dx.
  \]
  Uniform rotations are free, and long-wavelength twists are arbitrarily soft.

  Expanding about one vacuum gives one massive radial fluctuation and $m-1$
  transverse directions with no quadratic potential term—the variational
  analogue of Goldstone modes \cite{goldstone1962}.
\end{frame}

\begin{frame}{Connections encode relational comparison}
  \bookref{Chapter 3, Sections 8--9}
  A manifold-valued field does not by itself require a gauge field.
  As additional bundle/frame data, introduce
  \[
    A\in\Omega^1(\Omega,\mathfrak g),
    \qquad
    D=d+A,
  \]
  as an infinitesimal rule for comparing local orientations.

  Its curvature is
  \[
    F=dA+A\wedge A,
    \qquad
    [D_i,D_j]\phi=F_{ij}\cdot\phi.
  \]

  \begin{center}
    \begin{tikzpicture}[>=Latex,node distance=1.25cm]
      \node[draw,circle,fill=RelBlue!8] (x) {$x$};
      \node[draw,circle,fill=RelBlue!8,right=3.3cm of x] (y) {$y$};
      \draw[->,thick,bend left=28] (x) to node[above] {$\gamma_1$} (y);
      \draw[->,thick,bend right=28] (x) to node[below] {$\gamma_2$} (y);
      \node[right=0.5cm of y,align=left] {infinitesimal path discrepancy\\detected by curvature};
    \end{tikzpicture}
  \end{center}
\end{frame}

\begin{frame}{Curvature dynamics}
  \bookref{Chapter 3, Section 10}
  Given a positive-definite invariant inner product on $\mathfrak g$ and a background Hodge
  star, define
  \[
    S[A]=\int_\Omega\operatorname{Tr}(F\wedge *F).
  \]
  Compactly supported smooth variations of a smooth connection yield
  \[
    \delta F=D(\delta A),
    \qquad
    \boxed{D*F=0.}
  \]

  For $G=U(1)$ this gives $d*F=0$ (with $dF=0$); spacetime dynamics
  requires a Lorentzian metric;
  for nonabelian $G$ it has Yang--Mills form \cite{yangmills1954}.

  \begin{alertblock}{Interpretive status}
    Standard connection--curvature mathematics is being given a relational
    reading. The book does not derive Yang--Mills theory from the finite formula
    $I(S)=\log(n!/|\Aut(S)|)$.
  \end{alertblock}
\end{frame}

\begin{frame}{Where the continuum program stands}
  \bookref{Chapter 3, Sections 6--12}
  \begin{center}
    \begin{tikzpicture}[
      box/.style={draw,rounded corners,align=center,minimum width=2.25cm,
        minimum height=0.75cm},>=Latex,node distance=0.6cm]
      \node[box,fill=RelTeal!10] (done) {derived in explicit\\models};
      \node[box,fill=RelGold!12,right=of done] (interp) {geometric\\reinterpretation};
      \node[box,fill=RelRed!8,right=of interp] (open) {open unifying\\theory};
      \draw[->,thick] (done)--(interp);
      \draw[->,thick] (interp)--(open);
    \end{tikzpicture}
  \end{center}

  \begin{columns}[T]
    \column{0.32\textwidth}
      \small Block variational calculus; Hessian; Landau jet; $O(m)$ model.
    \column{0.32\textwidth}
      \small Local orientations as $G/H$ fields; connections as comparison;
      curvature as obstruction.
    \column{0.32\textwidth}
      \small Derive continuum geometry from microscopic ensembles; construct an
      RG flow; relate $I$ to Fisher--Rao geometry.
  \end{columns}
\end{frame}
